\section{Problem 1}

Let $\Omega = B_1^n(0)$ and $u(x)=|x|^{-s}$ with $s>0$.

\subsection{Part 1, Case 1: $u \in L^q(\Omega)$}
This must satisfy the finite q-integrability of the function $u(x)$, that is:
\[
\int_\Omega |u(x)|^q \, dx <\infty.
\]
We can determine the condition for the exponent $s$ from the above integration criteria in $L^q$ space. The only concern for this integration is that it will be extremely cumbersome to integrate in cartesian coordinate system. Now, the domain is given as $n$-dimensional unit ball ($\Omega = B_1^n(0)$). Thus, the accurate coordinate system is spherical polar coordinates. It will decompose the measure in radial measure ($r=|x|$) and angular measure ($\omega\in S^{n-1}$). The usual integration in $B_1^n(0)$ for any function $f(x)$ is in the following form:
\[
\int_{B_1^n(0)} f(x) \, dx = \int_{r=0}^1\int_{\omega\in S^{n-1}} f(r\omega)r^{n-1} \, d\sigma(\omega) \, dr.
\]
Applying this integration for $u(x)$ gives:
\begin{equation}
    \int_\Omega |u(x)|^q \, dx = \int_0^1\int_{S^{n-1}} |u(r\omega)|^q r^{n-1} \, d\sigma \, dr.
\end{equation}
Putting $u(x)=|x|^{-s}$ and $u(r\omega)=r^{-s}$ in (1), gives:
\[
\int_\Omega |x|^{-sq} \, dx = \int_0^1\int_{S^{n-1}} r^{-sq}r^{n-1} \, d\sigma \,dr.
\]
Decomposing the above equation gives:
\begin{equation}
    \int_\Omega |x|^{-sq} \, dx = \left(\int_{S^{n-1}} \, d\sigma \right)\int_0^1 r^{n-1-sq} \, dr.
\end{equation}
Now, the first integration gives surface measure on $n$-sphere:
\[
|S^{n-1}| = \frac{2\pi^{n/2}}{\Gamma(\frac{n}{2})},
\]
where $\Gamma(\cdot)$ is the usual gamma function. So (2) becomes:
\begin{equation}
    \int_\Omega |x|^{-sq} \, dx = \frac{2\pi^{n/2}}{\Gamma(\frac{n}{2})}\int_0^1 r^{n-1-sq} \, dr.
\end{equation}
The main objective here is that the integration expression of (3) must be finite for $u(x)\in L^q$. Now $\frac{2\pi^{n/2}}{\Gamma(\frac{n}{2})} < \infty$ as the domain is finite dimensional. Hence we need:
\[
\int_0^1 r^{n-1-sq} \, dr <\infty.
\]
This will converge when:
\[
n-1-sq > -1.
\]
So the condition of $s$ for $u(x)\in L^q(\Omega)$ is:
\[
\boxed{
n > sq.
}
\]

\subsection{Case 2: $u\in W^{1,p}(\Omega)$}

From the definition of Sobolev space:
\[
u\in W^{1,p} \iff u\in L^p(\Omega), \, \nabla u\in L^p(\Omega).
\]

\subsection*{Case 2(a): $u\in L^p(\Omega)$}
It will give the similar condition as $u\in L^q(\Omega)$, with $q$ replaced by $p$:
\[
\boxed{
n>sp.
}
\]

\subsection*{Case 2(b): $\nabla u\in L^p(\Omega)$}
The weak derivative expression of $u(x)$ is:
\[
\nabla u = -s|x|^{-s-2}x.
\]
So, the absolute value of the derivative is:
\begin{equation}
    |\nabla u(x)| = s|x|^{-s-1}.
\end{equation}
$\nabla u\in L^p(\Omega) \implies$:
\[
\int_\Omega |\nabla u|^p \, dx = \frac{2\pi^{n/2}s}{\Gamma(\frac{n}{2})} \int_0^1 r^{-p(s+1)}r^{n-1} \, dr < \infty.
\]
That is:
\[
\int_\Omega |\nabla u|^p \, dx = \frac{2\pi^{n/2}s}{\Gamma(\frac{n}{2})} \int_0^1 r^{n-1-p(s+1)} \, dr.
\]
This $p$-integrability will be finite only if:
\[
\int_0^1 r^{n-1-p(s+1)} \, dr <\infty.
\]
This integration will converge if:
\[
n-1-p(s+1) > -1.
\]
Which gives:
\[
s < \frac{n}{p}-1.
\]
Hence, the condition of $s$ for $u \in W^{1,p}(\Omega)$ is:
\[
\boxed{
n> sp \quad \text{and} \quad s<\frac{n}{p}-1.
}
\]

Since $u\in C^1(\Omega\setminus\{0\})$, the singular set $\{0\}$ has measure zero, and both $u$ and its classical gradient belong to $L^p(\Omega)$, it follows that the classical gradient coincides with the weak gradient. Hence, $\nabla u$ is the weak derivative and $u\in W^{1,p}(\Omega)$.

\subsection{Part 2: Sharp Sobolev exponent $p^\ast$}
For $1\leq p < n$, Sobolev embedding theorem gives that:
\[
W^{1,p}(\Omega) \hookrightarrow L^{p^\ast}(\Omega),
\]
where $p^\ast=\frac{np}{n-p}$ and satisfies the Sobolev inequality:
\[
\|u\|_{L^{p^\ast}} \leq C\|u\|_{W^{1,p}}.
\]
Sharp Sobolev exponent means $p^\ast$ cannot be replaced by any larger exponent than $p^\ast$. This implies:
\[
W^{1,p}(\Omega) \not\hookrightarrow L^q(\Omega)
\]
for any $q>p^\ast$. 
From the previous finding condition for $u\in W^{1,p}(\Omega)$:
\[
s<\frac{n}{p}-1.
\]
Taking $\epsilon >0$ gives:
\begin{equation}
    s=\frac{n}{p}-1-\epsilon.
\end{equation}
Putting (5) in the inequality $sq<n$ (from the condition of $s$ for $u\in L^q(\Omega)$):
\[
\left(\frac{n}{p}-1-\epsilon \right)q < n.
\]
Taking $\epsilon \to 0$ gives:
\[
q \leq \frac{np}{n-p}=p^\ast.
\]
This proves that the Sobolev exponent $p^\ast$ is sharp.

\section{Problem 2}

It is given that the dimension $n\geq 2$. The function $u(x)=\frac{x}{|x|}\in B_1^n(0)$. We need to find the range of $p$ such that $u\in W^{1,p}(B_1^n(0))$. From the definition of Sobolev space:
\[
u\in W^{1,p}(B_1^n(0)) \iff u\in L^p(B_1^n(0)), \, \nabla u\in L^p(B_1^n(0)).
\]

\subsection{Case 1: $u\in L^p(B_1^n(0))$}
$u(x)=\frac{x}{|x|}$ for all $x\in B_1^n(0)\setminus\{0\}$:
\[
|u(x)|=\left|\frac{x}{|x|}\right|=\frac{|x|}{|x|}=1 \quad \text{for all } x\neq 0.
\]
Hence, for any $p\geq 1$:
\[
|u(x)|^p= 1.
\]
Thus, the p-th integration in spherical polar coordinates is:
\[
\int_{B_1^n(0)}|u(x)|^p \, dx = \int_{B_1^n(0)} \, dx=|B_1^n(0)|.
\]
Now, as the unit ball is in finite dimension, it will always be bounded. This says:
\[
\int_{B_1^n(0)}|u(x)|^p \, dx < \infty \quad \text{for every } 1\leq p\leq\infty.
\]

\subsection{Case 2: $\nabla u\in L^p(B_1^n(0))$}
For $u(x)=\frac{x}{|x|}$, the absolute value of the weak derivative is:
\[
|\nabla u(x)| = \frac{1}{|x|}.
\]
Now, the $L^p$ integrability in spherical polar coordinates gives:
\[
\int_{B_1^n(0)}|\nabla u(x)|^p \, dx= S_M\int_0^1 r^{n-1-p} \, dr,
\]
where $S_M= \frac{2\pi^{n/2}}{\Gamma(\frac{n}{2})}$ is the surface measure, which is finite. This integrability will converge $\iff$:
\[
n-1-p>-1.
\]
\[
\iff p<n.
\]
Hence the range of $p$ for $u\in W^{1,p}(B_1^n(0))$ is:
\[
\boxed{
1\leq p < n.
}
\]

Since $u\in C^1(\Omega\setminus\{0\})$, the singular set $\{0\}$ has measure zero, and both $u$ and its classical gradient belong to $L^p(\Omega)$, it follows that the classical gradient coincides with the weak gradient. Hence, $\nabla u$ is the weak derivative and $u\in W^{1,p}(B_{1}^n(0))$.

\subsection{Part 2: $n=1$}
\begin{proof}
For $n=1$, take the domain as an open interval $(a,b)$. Let $u\in W^{1,p}(a,b)$. Then, for almost every $x, y\in (a,b)$, by the fundamental theorem of calculus:
\[
u(x)-u(y)=\int_x^y u'(t)\, dt.
\tag{1}
\]
Taking absolute value on (1) and using Jensen's inequality:
\[
|u(x)-u(y)|=\left|\int_x^y u'(t)\, dt\right|\le \int_y^x |u'(t)|\, dt.
\]
Using H\"older:
\[
|u(x)-u(y)| \le |x-y|^{1-1/p}\|u'\|_{L^p(a,b)}.
\]
Hence:
\[
|u(x)-u(y)| \le C|x-y|^\alpha \quad \text{with } \alpha=1-\frac{1}{p}>0.
\]
This is the H\"older estimate. Hence:
\[
u\in C^{0,\alpha}([a,b])\subset C([a,b]).
\]
Thus, for $n=1$, every Sobolev function is continuous. The above proof showed $u$ is H\"older continuous (for $p>1$) or absolutely continuous (for $p=1$).
\end{proof}

\section{Problem 3}
\begin{proof}
The given function is:
\[
u(x) = |\log |x||^s.
\]
The domain is $B_{1/2}^2(0)$. Hence the dimension is $n=2$. From the definition of Sobolev function:
\[
u(x)\in W^{1,2}(B_{1/2}^2(0)) \iff u\in L^2(B_{1/2}^2(0))\, ,\, \nabla u(x)\in L^2(B_{1/2}^2(0)).
\]
So, this problem is strictly about $p=2$ for 2D. 

\subsection{Case 1: $u(x)\in L^2(B_{1/2}^2(0))$}
This is a radial function:
\[
u(x)=|\log r|^s.
\]
The condition for square integrability is:
\[
\int_{B_{1/2}^2(0)} |u|^2 \, dx < \infty.
\]
The integration of the radial function in 2D spherical polar coordinates is:
\begin{equation}
    \int_{B_{1/2}^2(0)} |u|^2 \, dx = \int_0^{1/2} \int_{S^1} |\log r|^{2s} \, r \, d\sigma \, dr.
\end{equation}
The surface measure gives $|S^1|=2\pi$. Hence, (6) becomes:
\begin{equation}
    \int_{B_{1/2}^2(0)} |u|^2 \, dx = 2\pi\int_0^{1/2} |\log r|^{2s} \, r \, dr.
\end{equation}
Now, make the change of variables:
\begin{align*}
    t &= -\log r, \\
    r &= e^{-t}, \\
    dr &= -e^{-t}\, dt.
\end{align*}
The lower integration limit becomes $\log 2$ and upper integration limit is $\infty$. Hence (7) becomes:
\[
\int_{B_{1/2}^2(0)} |u|^2 \, dx = 2\pi\int_{\log 2}^\infty t^{2s}e^{-2t} \, dt.
\]
Make another change of variables: $z=2t$, so $t=z/2$, $dt=dz/2$:
\[
\int_{B_{1/2}^2(0)} |u|^2 \, dx = \frac{2\pi}{2^{2s+1}}\int_{2\log 2}^\infty z^{2s} e^{-z} \, dz.
\]
That is:
\begin{equation}
    \int_{B_{1/2}^2(0)} |u|^2 \, dx = \frac{\pi}{2^{2s}}\Gamma(2s+1, 2\log 2),
\end{equation}
where $\Gamma(\cdot,\cdot)$ is the upper incomplete gamma function. The prefactor is finite, and $\Gamma(2s+1, 2\log 2)$ is finite for all $s>-1/2$. Since $0<s<1$, the integral is finite:
\[
\int_{B_{1/2}^2(0)} |u|^2 \, dx < \infty.
\]
This proves:
\[
\boxed{u(x) \in L^2(B_{1/2}^2(0)).}
\]

\subsection{Case 2: $\nabla u\in L^2(B_{1/2}^2(0))$}
The derivative of the function is:
\[
|\nabla u(x)| = s|\log r|^{s-1}\frac{1}{r}.
\]
Now, test square-integrability:
\[
\int_{B_{1/2}^2(0)} |\nabla u(x)|^2 \, dx = \int_0^{1/2}\int_{S^1} s^2|\log r|^{2s-2}\frac{1}{r^2}r \, d\sigma\, dr.
\]
With further simplification:
\begin{equation}
    \int_{B_{1/2}^2(0)} |\nabla u(x)|^2 \, dx = 2\pi s^2\int_0^{1/2} |\log r|^{2s-2}\frac{1}{r} \, dr.
\end{equation}
Make the change of variables:
\begin{align*}
    t &= -\log r, \\
    \frac{1}{r}\, dr &= -dt.
\end{align*}
This gives lower limit $\log 2$ and upper limit $\infty$:
\[
\int_{B_{1/2}^2(0)} |\nabla u(x)|^2 \, dx = 2\pi s^2\int_{\log 2}^\infty t^{2s-2} \, dt.
\]
This integral converges if and only if:
\[
2s-2 < -1 \iff s < \frac{1}{2}.
\]
The case $s=0$ gives $u(x)=1$, which is constant and hence $u\in W^{1,2}$ (trivial case). 

Thus:
\[
\boxed{
u\in W^{1,2}(B_{1/2}^2(0)) \quad \text{for all } 0\le s<1/2.
}
\]

Since $u\in C^1(B_{1/2}^2(0)\setminus\{0\})$ and both $u$ and $\nabla u$ belong to $L^2(B_{1/2}^2(0))$, the classical gradient coincides with the weak gradient. The singular set $\{0\}$ has measure zero.
\end{proof}

\section{Problem 4}
\begin{proof}
Let $n \ge 2$ and let $U = B_1^n(0) \subset \mathbb{R}^n$. Define
\[
u(x) = \log\!\big(\log(1 + \tfrac{1}{|x|})\big).
\]

We shall show that $u \in W^{1,n}(U)$, i.e.,
\[
u \in L^n(U) \quad \text{and} \quad \nabla u \in L^n(U).
\]

\noindent
\subsection*{Case 1: $u \in L^n(U)$}
Since $u$ is smooth away from $x=0$, the only possible singularity is at the origin. As $r = |x| \to 0$,
\[
u(x) = \log(\log(1 + 1/r)) \sim \log(\log(1/r)),
\]
which grows slower than any power of $r^{-1}$. Hence
\[
\int_{B_1(0)} |u(x)|^n \, dx < \infty,
\]
and therefore $u \in L^n(U)$.

\noindent
\subsection*{Case 2: $\nabla u \in L^n(U)$}
Let $r = |x|$ and define
\[
f(r) = \log(\log(1 + 1/r)).
\]
Then $u(x) = f(|x|)$ and by the chain rule,
\[
f'(r)
= \frac{1}{\log(1+1/r)} \cdot \frac{1}{1+1/r} \cdot \left(-\frac{1}{r^2}\right)
= -\frac{1}{r(r+1)\log(1+1/r)}.
\]
Since
\[
\nabla u(x) = f'(r)\frac{x}{r},
\]
we obtain
\[
|\nabla u(x)| = |f'(r)| 
= \frac{1}{r(r+1)\log(1+1/r)}.
\]

We compute:
\[
\int_{B_1^n(0)} |\nabla u(x)|^n \, dx
= C_n \int_0^1 
\frac{1}{r^n(r+1)^n(\log(1+1/r))^n}
r^{n-1} \, dr
\]
\[
= C_n \int_0^1 
\frac{1}{r(\log(1+1/r))^n} \cdot \frac{1}{(r+1)^n} \, dr.
\]
Since $0 < r < 1$, we have $(r+1)^{-n} \le 1$, hence
\[
\int_{B_1^n(0)} |\nabla u|^n dx
\le C_n \int_0^1 \frac{dr}{r(\log(1+1/r))^n}.
\]

Set
\[
t = \log(1 + 1/r).
\]
Then
\[
dt = -\frac{1}{r(r+1)} \, dr,
\quad \text{hence} \quad
\frac{dr}{r} = -(r+1)\, dt \le 2\, dt.
\]
Changing variables and reversing the limits of integration:
\begin{align*}
    \int_0^1\frac{dr}{r(\log(1+1/r))^n}
    &=\int_{\infty}^{\log 2}\frac{-(r+1)}{t^n}\,dt \\
    &=\int_{\log 2}^{\infty}\frac{r+1}{t^n}\, dt.
\end{align*}
Since $0<r<1$ implies $1<r+1<2$, we have:
\[
\int_0^1 \frac{dr}{r(\log(1+1/r))^n}
\le 2 \int_{\log 2}^{\infty} \frac{dt}{t^n}.
\]
The integral on the right converges if and only if $n>1$, which holds by assumption (indeed $n\ge 2$).

Thus,
\[
\int_{B_1^n(0)} |\nabla u|^n dx < \infty.
\]

We have shown that
\[
u \in L^n(B_1^n(0)) 
\quad \text{and} \quad 
\nabla u \in L^n(B_1^n(0)).
\]
Hence
\[
u \in W^{1,n}(B_1^n(0)).
\]
\end{proof}

\section{Problem 5}
We have to prove the following theorem:
\begin{theorem}
    For every $1\le p<\infty$, there exists a constant $c=c(n, p) >0$ such that
    \[
    \fint_{B_R(x)} \left|u - (u)_{B_R(x)}\right |^p \le cR^p \fint_{B_R(x)} |\nabla u|^p
    \]
    for any ball $B_R(x) \subset \mathbb{R}^n$ and every $u\in W^{1,p}(\mathbb{R}^n)$.
\end{theorem}

\begin{proof}
Here, $(u)_{B_R(x)}$ represents the average or mean integral of $u$ over $B_R(x)$:
\[
(u)_{B_R(x)} :=\frac{1}{|B_R(x)|}\int_{B_R(x)} u(y) \, dy
= \fint_{B_R(x)} u(y) \, dy.
\]

Initially, assume $u\in C^1(\mathbb{R}^n)$. Then:
\begin{equation}
\fint_{B_R(x)} |u(y) - (u)_{B_R(x)}|^p \, dy = \fint_{B_R(x)} \left|\fint_{B_R(x)} (u(y) - u(w)) \, dw\right|^p \, dy.
\end{equation}
This is valid because:
\[
\fint_{B_R(x)} (u(y)-u(w)) \, dw = u(y) - \fint_{B_R(x)} u(w) \, dw = u(y) -(u)_{B_R(x)}.
\]

Applying Jensen's inequality to the right side of (9):
\[
\fint_{B_R(x)} \left|\fint_{B_R(x)} (u(y) - u(w)) \, dw\right|^p \, dy \le \fint_{B_R(x)} \fint_{B_R(x)} |u(y) -u(w)|^p \, dy \, dw.
\]
Thus:
\begin{equation}
    \fint_{B_R(x)} |u(y) - (u)_{B_R(x)}|^p \, dy \le \fint \fint_{B_R(x)} |u(y) -u(w)|^p \, dy \, dw.
\end{equation}

We use the following lemma:
\begin{lemma}
    For every $1\le p < \infty$, there exists $c>0$ depending only on $n$ and $p$ such that 
    \[
    \int_{B_R(x)} |u(y)-u(w)|^p \, dy \le cR^{p-1}\int_{B_R(x)} \frac{|\nabla u(y)|^p}{|y-w|^{n-1}}\, dy
    \]
    for every ball $B_R(x) \subset \mathbb{R}^n$, every $w\in B_R(x)$, and every $u\in W^{1,p}(\mathbb{R}^n)$.
\end{lemma}

Using this lemma to estimate the right side of (10):
\begin{align}
\fint \fint_{B_R(x)} |u(y) -u(w)|^p \, dy \, dw
&\le c\fint_{B_R(x)} R^{p-1} \int_{B_R(x)} \frac{|\nabla u(w)|^p}{|y-w|^{n-1}}\, dw \, dy.
\end{align}
There is no $n$ in the exponent of $R$ in (11) due to the mean integral. Since $|y-w| \le 2R$:
\[
\int_{B_R(x)} \frac{|\nabla u(w)|^p}{|y-w|^{n-1}}\, dw \, dy \le C R \int_{B_R(x)} |\nabla u(w)|^p \, dw.
\]
Hence:
\begin{align}
    \fint_{B_R(x)} |u(y) - (u)_{B_R(x)}|^p \, dy
    &\le c R^{p-1} \fint_{B_R(x)} R \int_{B_R(x)} |\nabla u(w)|^p \, dw \\
    &= c R^p \fint_{B_R(x)} |\nabla u(w)|^p \, dw.
\end{align}
Thus:
\[
\boxed{
\fint_{B_R(x)} \left|u - (u)_{B_R(x)}\right |^p \le cR^p \fint_{B_R(x)} |\nabla u|^p.
}
\]
The result extends to $u\in W^{1,p}(\mathbb{R}^n)$ by density.
\end{proof}

\section{Problem 6}
\begin{theorem}[Poincar\'e inequality on convex domains]
Let $\Omega \subset \mathbb{R}^n$ be open, bounded, convex and $C^1$.  
Then for every $1 \le p < \infty$ there exists a constant 
$c = c(n,p,\Omega) > 0$ such that
\[
\fint_\Omega |u-(u)_\Omega|^p \, dx 
\le c \fint_\Omega |\nabla u|^p \, dx
\quad \text{for all } u \in W^{1,p}(\Omega).
\]
\end{theorem}

\begin{proof}
Fix $u \in W^{1,p}(\Omega)$. Since $\Omega$ is convex, for every $x,y \in \Omega$ the line segment $[x,y] \subset \Omega$, and by the fundamental theorem of calculus,
\[
u(x)-u(y)= \int_0^1 \nabla u(x+t(y-x)) \cdot (y-x)\, dt.
\]
Taking absolute values:
\[
|u(x)-u(y)|
\le |x-y| \int_0^1 |\nabla u(x+t(y-x))| \, dt.
\tag{13}
\]
Using the identity
\[
u(x)-(u)_\Omega = \fint_\Omega (u(x)-u(y))\, dy,
\]
and applying (13), we obtain:
\[
|u(x)-(u)_\Omega|
\le \fint_\Omega |x-y| \int_0^1 |\nabla u(x+t(y-x))| \, dt \, dy.
\]
Interchanging the integrals:
\[
|u(x)-(u)_\Omega|
\le \int_0^1 \fint_\Omega |x-y|\,|\nabla u(x+t(y-x))| \, dy \, dt.
\tag{14}
\]
Since $|x-y| \le \operatorname{diam}(\Omega)$, (14) becomes:
\[
|u(x)-(u)_\Omega|
\le \operatorname{diam}(\Omega)
\int_0^1 \fint_\Omega |\nabla u(x+t(y-x))| \, dy \, dt.
\tag{15}
\]
Raising to the power $p$, integrating in $x$, and applying Jensen's inequality:
\[
\int_\Omega |u-(u)_\Omega|^p \, dx
\le C(\Omega)
\int_\Omega \int_0^1 \fint_\Omega 
|\nabla u(x+t(y-x))|^p \, dy \, dt \, dx,
\tag{16}
\]
where $C(\Omega) = \operatorname{diam}(\Omega)^p$.

For fixed $t \in (0,1)$ and $y \in \Omega$, the map
\[
x \mapsto z = x+t(y-x)
\]
maps $\Omega$ into itself and is Lipschitz with Jacobian bounded independently of $x,y,t$. Hence:
\[
\int_\Omega |\nabla u(x+t(y-x))|^p \, dx
\le C(n) \int_\Omega |\nabla u(z)|^p \, dz.
\]
Substituting into (16):
\[
\int_\Omega |u-(u)_\Omega|^p \, dx
\le C(\Omega)C(n)
\int_0^1 \int_\Omega |\nabla u(z)|^p \, dz \, dt.
\]
Since $\int_0^1 dt = 1$, we conclude:
\[
\int_\Omega |u-(u)_\Omega|^p \, dx
\le C(\Omega)C(n)
\int_\Omega |\nabla u|^p \, dx.
\]
Dividing by $|\Omega|$:
\[
\fint_\Omega |u-(u)_\Omega|^p \, dx
\le c(n,p,\Omega)\, \fint_\Omega |\nabla u|^p \, dx,
\]
where
\[
c(n,p,\Omega) = C(n)\operatorname{diam}(\Omega)^p.
\]
This completes the proof.
\end{proof}

\section{Problem 7}
\begin{theorem}
    Let $\Omega\subset \mathbb{R}^n$ be open, bounded and $C^1$. For every $1\le p< \infty$, there exists $c=c(n, p, \Omega)>0$ such that 
    \[
    \fint_\Omega |u-(u)_\Omega|^p \le c\fint_\Omega |\nabla u|^p
    \]
    for every $u\in W^{1,p}(\Omega)$.
\end{theorem}

\begin{proof}
We prove by contradiction. Assume no constant $c>0$ exists. Then for every integer $k\ge 1$, choose $u_k\in W^{1,p}(\Omega)$ with:
\[
\int_\Omega |u_k-(u_k)_\Omega|^p > k\int_\Omega |\nabla u_k|^p.
\tag{17}
\]
Define $w_k= u_k-(u_k)_\Omega$. Then:
\[
\int_\Omega w_k = 0,
\]
so (17) becomes:
\[
\int_\Omega |w_k|^p > k\int_\Omega |\nabla w_k|^p.
\tag{18}
\]
Normalize:
\[
v_k = \frac{w_k}{\|w_k\|_{L^p(\Omega)}}.
\]
Then:
\[
\|v_k\|_{L^p(\Omega)} = 1, \qquad \int_\Omega v_k = 0.
\]
Dividing (18) by $\|w_k\|_{L^p(\Omega)}^p$:
\[
1 > k \int_\Omega |\nabla v_k|^p.
\]
Thus:
\[
\int_\Omega |\nabla v_k|^p < \frac{1}{k}.
\]
So $\|\nabla v_k\|_{L^p(\Omega)} \to 0$ as $k\to\infty$. Also:
\[
\|v_k\|_{W^{1,p}(\Omega)}^p = \|v_k\|_{L^p(\Omega)}^p + \|\nabla v_k\|_{L^p(\Omega)}^p \le 1+\frac{1}{k} \le 2.
\]
Hence $\{v_k\}$ is bounded in $W^{1,p}(\Omega)$.

Since $\Omega$ is bounded and $C^1$, the Rellich-Kondrachov compactness theorem implies that $W^{1,p}(\Omega)$ is compactly embedded in $L^p(\Omega)$. Therefore, there exists a subsequence and a function $v\in L^p(\Omega)$ such that $v_k\to v$ strongly in $L^p(\Omega)$. In particular:
\[
\|v\|_{L^p(\Omega)} = \lim_{k\to\infty}\|v_k\|_{L^p(\Omega)}=1,
\]
and because $\int_\Omega v_k = 0$ for all $k$, we also have $\int_\Omega v = 0$.

Moreover, for any test function $\phi\in C_c^\infty(\Omega)$ and any $i=1,\dots,n$,
\[
\int_\Omega v\,\partial_i\phi\, dx
=\lim_{k\to\infty}\int_\Omega v_k\,\partial_i\phi\, dx
=-\lim_{k\to\infty}\int_\Omega \partial_i v_k\,\phi\, dx.
\]
Since $\|\partial_i v_k\|_{L^p(\Omega)} \le \|\nabla v_k\|_{L^p(\Omega)} \to 0$, the right-hand side tends to zero. Hence:
\[
\int_\Omega v\,\partial_i\phi\, dx = 0 \quad \text{for all } \phi\in C_c^\infty(\Omega).
\]
This means $v$ has weak derivative zero. Therefore $v$ is constant on each connected component of $\Omega$. Since $\Omega$ is connected, $v$ is constant everywhere. Together with $\int_\Omega v = 0$, this forces $v=0$ almost everywhere, contradicting $\|v\|_{L^p(\Omega)}=1$.

Hence the assumption was false and there exists $c=c(n,p,\Omega)>0$ such that:
\[
\fint_\Omega |u-(u)_\Omega|^p \le c\fint_\Omega |\nabla u|^p.
\]
This completes the proof.
\end{proof}

\section{Problem 8}
\begin{proof}
As $W^{1,\infty}(\Omega)\subset W^{1,p}(\Omega)$ for any $n<p<\infty$, according to the Sobolev embedding theorem:
\begin{theorem}
    Let $\Omega\subset\mathbb{R}^n$ be open, bounded and smooth and let $n<p<\infty$. Then $W^{1,p}(\Omega)$ continuously embeds into $C^{0,\tau}(\overline{\Omega})$ for every $0\le\tau\le 1-\frac{n}{p}$.
\end{theorem}
Hence, for any $x,y\in\overline{\Omega}$:
\[
|u(x)-u(y)|\le c|x-y|^{1-\frac{n}{p}}\|u\|_{W^{1,p}(\Omega)}.
\]
Letting $p\to\infty$ and noting that:
\[
\lim_{p\to\infty}\|u\|_{W^{1,p}(\Omega)}=\|u\|_{W^{1,\infty}(\Omega)},
\]
we obtain:
\[
[u]_{C^{0,1}(\overline{\Omega})} := \sup_{x,y\in \overline{\Omega}}\frac{|u(x)-u(y)|}{|x-y|}\le c\|u\|_{W^{1,\infty}(\Omega)}.
\]
So we have established the continuous embedding:
\[
W^{1,\infty}(\Omega) \subset C^{0,1}(\overline{\Omega}).
\tag{19}
\]

Now prove the reverse inclusion. By Rademacher's theorem, every Lipschitz function $u:\Omega\to\mathbb{R}$ is differentiable almost everywhere, and 
\[
|\nabla u(x)|\le [u]_{C^{0,1}} \quad \text{for almost every } x.
\]
Hence $\partial_i u\in L^\infty(\Omega)$ and $u\in L^\infty(\Omega)$. Since $u$ is a.e. differentiable and bounded, integration by parts is justified:
\[
\int_\Omega u\, \partial_i\phi = -\int_\Omega \partial_i u\,\phi \quad \forall \phi \in C_c^\infty(\Omega).
\]
Thus the classical derivative equals the weak derivative. So:
\[
u\in L^\infty(\Omega), \qquad \nabla u\in L^\infty(\Omega).
\]
Therefore $u\in W^{1,\infty}(\Omega)$ and:
\[
\|u\|_{W^{1,\infty}}\le \|u\|_{L^\infty} + [u]_{C^{0,1}}.
\]
Hence:
\[
C^{0,1}(\overline{\Omega})\subset W^{1,\infty}(\Omega),
\tag{20}
\]
with equivalent norms:
\[
\|u\|_{W^{1,\infty}}\cong \|u\|_{L^\infty}+[u]_{C^{0,1}}.
\]
Combining (19) and (20):
\[
\boxed{
W^{1,\infty}(\Omega) = C^{0,1}(\overline{\Omega}).
}
\]
\end{proof}

\section{Problem 9}
\begin{proof}
Let $A\subset\mathbb{R}^n$ be measurable with $0<|A|<\infty$. Let $f\in L^p(A)$, $(f)_A=\frac{1}{|A|}\int_A f$, and $\lambda\in\mathbb{R}$ arbitrary.

First decompose:
\[
f-(f)_A = (f-\lambda)+(\lambda-(f)_A).
\tag{21}
\]
Using Young's inequality $|a+b|^p\le 2^{p-1}(|a|^p+|b|^p)$ for $p\ge 1$:
\[
|f-(f)_A|^p\le 2^{p-1}(|f-\lambda|^p+|\lambda-(f)_A|^p).
\tag{22}
\]
Integrating (22) over $A$:
\[
\int_A|f-(f)_A|^p \le 2^{p-1}\left(\int_A |f-\lambda|^p + |A|\,|\lambda-(f)_A|^p\right).
\tag{23}
\]
Applying Jensen's inequality:
\[
|\lambda-(f)_A|=\left|\frac{1}{|A|}\int_A(\lambda-f(x))\, dx\right|\le \frac{1}{|A|}\int_A |f-\lambda|\, dx.
\tag{24}
\]
Applying H\"older's inequality to (24):
\[
\frac{1}{|A|}\int_A |f-\lambda| \le\left(\frac{1}{|A|}\int_A|f-\lambda|^p\right)^{1/p}.
\tag{25}
\]
From (24) and (25):
\[
|\lambda-(f)_A|^p \le \frac{1}{|A|}\int_A|f-\lambda|^p.
\tag{26}
\]
Multiplying (26) by $|A|$:
\[
|A|\,|\lambda-(f)_A|^p \le \int_A|f-\lambda|^p.
\tag{27}
\]
Substituting (27) into (23):
\[
\int_A|f-(f)_A|^p\le 2^{p-1}\left(\int_A |f-\lambda|^p+\int_A|f-\lambda|^p\right)
=2^p\int_A |f-\lambda|^p.
\]
Thus:
\[
\boxed{
\int_A|f-(f)_A|^p\le 2^p\int_A |f-\lambda|^p.
}
\]
Dividing by $|A|$ gives the averaged form. This proves the desired inequality.
\end{proof}

\end{document}